\chapter{Preface}
\markboth{Preface}{Preface}

\section*{Why this book exists}

Imagine a small fleet of drones that must fly through the same airspace at the same
time. Each drone has a job to do and a route to fly. Before take-off, a planner must
make sure that no two routes collide. During the flight, something unplanned happens:
another drone, not part of the fleet and not talking to it, crosses the formation. The
fleet must notice it, guess where it is going, get out of its way, and then return to
the plan---or, if that is impossible, make a new plan and check it against the rest of
the fleet.

Every part of that story is a well-studied algorithmic problem. Planning a route is
\emph{graph search}. Planning routes for many drones at once without collisions is
\emph{multi-agent path finding}. Reacting to a surprise in a fraction of a second is
\emph{local collision avoidance}. Guessing where the intruder is going is
\emph{tracking and trajectory prediction}. Making a new plan without starting from
scratch is \emph{incremental replanning}. Keeping the fleet in formation and within
radio range while doing all of this is \emph{control}. The individual algorithms are
established and published. What is not settled is how to combine them into one system
that keeps the guarantees of the global plan, reacts fast enough to be safe, and can be
evaluated honestly.

This book teaches all of those algorithms, one at a time, from first principles, and
then shows how they fit together into a hybrid swarm planner. It was written to
accompany a twelve-week research training plan; the plan is reproduced, with
week-by-week reading and coding assignments, in \cref{ch:appA}.

\section*{Who should read it}

The book is for a graduate student or engineer who is starting research on drone
swarms, or on multi-robot motion planning more generally, and who wants to
\emph{implement} the algorithms rather than merely recognise their names. You should be
comfortable with Python and with the mathematics of a first course in linear algebra,
calculus and probability. Everything else is introduced when it is needed, and
\cref{ch:appB} collects the mathematical facts that the book relies on. No prior
knowledge of robotics or planning is assumed.

\section*{What you will be able to do}

When you have worked through the book and its exercises you will be able to
\begin{itemize}
  \item implement and explain \astar and space-time \astar, and say exactly why they
        return optimal paths under the required conditions;
  \item repair a path with \dstarlite instead of replanning from scratch, and measure
        when that is worth it;
  \item represent and detect the conflicts of a multi-agent plan, and solve the
        problem optimally with conflict-based search or approximately, with a
        guaranteed bound, with \ecbs;
  \item explain velocity obstacles, reciprocal velocity obstacles and \orca using
        relative position and velocity, and use them to keep drones apart;
  \item track an unknown drone from noisy observations with a Kalman filter, and
        compare a constant-velocity predictor with a learned one;
  \item trigger local avoidance from a time-to-collision or safety-horizon rule, return
        to the global route, or replan;
  \item enforce formation and communication-range constraints, and
  \item evaluate a hybrid planner with reproducible metrics and controlled scenarios.
\end{itemize}

\section*{How the book is organised}

\begin{description}[style=nextline,leftmargin=1.5em]
  \item[Part I, Foundations.] \Cref{ch:ch01} states the problem and previews the
        hybrid architecture that the whole book builds towards. \Cref{ch:ch02} is a
        toolbox: graphs, grids, configuration space, time-indexed paths, simple motion
        models, the geometry of moving discs, Gaussians and priority queues.
  \item[Part II, Single-Agent Graph Search.] Dijkstra's algorithm (\cref{ch:ch03}),
        \astar with admissible and consistent heuristics and its space-time version
        (\cref{ch:ch04}), incremental replanning with \lpastar and \dstarlite
        (\cref{ch:ch05}) and anytime search with \arastar (\cref{ch:ch06}).
  \item[Part III, Multi-Agent Path Finding.] The problem and its conflicts
        (\cref{ch:ch07}), prioritized planning (\cref{ch:ch08}), conflict-based
        search (\cref{ch:ch09}), bounded-suboptimal search with \ecbs
        (\cref{ch:ch10}), and the coupled and constructive alternatives \mstar,
        Push-and-Swap and Push-and-Rotate (\cref{ch:ch11}).
  \item[Part IV, Local and Reactive Collision Avoidance.] Velocity obstacles
        (\cref{ch:ch12}), reciprocal avoidance with RVO and \orca (\cref{ch:ch13}),
        the dynamic window approach (\cref{ch:ch14}) and artificial potential fields
        (\cref{ch:ch15}).
  \item[Part V, Sampling-Based Motion Planning.] \rrt (\cref{ch:ch16}) and the
        asymptotically optimal \rrtstar with informed sampling (\cref{ch:ch17}).
  \item[Part VI, Tracking and Prediction.] The Kalman filter (\cref{ch:ch18}), its
        nonlinear relatives and the particle filter (\cref{ch:ch19}), and trajectory
        prediction from constant-velocity baselines to LSTMs and Transformers
        (\cref{ch:ch20}).
  \item[Part VII, Optimization and Control.] Model predictive control
        (\cref{ch:ch21}), mixed-integer linear programming (\cref{ch:ch22}) and
        consensus and formation control (\cref{ch:ch23}).
  \item[Part VIII, Putting It Together.] The hybrid collision-avoidance architecture
        (\cref{ch:ch24}) and how to evaluate it (\cref{ch:ch25}).
\end{description}

\section*{How to read it}

Read the chapters in order if you can: each part builds on the previous ones, and the
notation is introduced once. If you are short of time, the training plan's reading
order is a good fast path: \cref{ch:ch04} (\astar), \cref{ch:ch05} (\dstarlite),
\cref{ch:ch09} (\cbs), \cref{ch:ch07} (the MAPF problem), \cref{ch:ch12,ch:ch13}
(velocity obstacles and \orca), \cref{ch:ch14} (DWA), \cref{ch:ch16,ch:ch17}
(\rrt and \rrtstar), \cref{ch:ch21} (MPC) and \cref{ch:ch20} (trajectory prediction),
then \cref{ch:ch24}. Every chapter ends with a box that says where its algorithm sits
in the hybrid system, so you can also read by \emph{layer}: the global planner
(Parts II and III), the local safety layer (Part IV), the prediction layer (Part VI)
and the constraints (Part VII).

\section*{Conventions}

Every chapter opens with a box of \emph{learning objectives} and closes with a
\emph{summary} box and exercises. Inside a chapter you will meet five kinds of coloured
boxes: \emph{key idea} (green) states the one thing to remember; \emph{in the drone
system} (purple) says how the algorithm is used in the hybrid architecture;
\emph{pitfall} (orange) warns about a mistake that is easy to make; \emph{historical
note} (grey) gives context; and plain notes carry their own titles. Definitions,
examples, theorems and remarks share one counter per chapter, so ``Example~4.3'' comes
after ``Definition~4.2''. Exercises are marked with one to three stars
(\difficulty{1} routine, \difficulty{2} needs thought, \difficulty{3} a small
project), and hints or solutions for a selection of them are in \cref{ch:appC}. Where
the training plan assigns a priority to an algorithm it is shown as a tag such as
\priority{Essential}. Pseudocode is written in a uniform style with numbered lines;
the text refers to line numbers. Vectors are bold ($\pos$, $\vel$), matrices are bold
capitals ($\mat{A}$), and the full notation is listed in the table that follows this
preface.

\section*{The code}

Every algorithm in the book comes with a \python{} implementation in the
\code{code/} folder of the \LaTeX{} project. The implementations are deliberately
simple---plain \python{} and NumPy, a few hundred lines each, no framework---because
the goal is that you can read them in one sitting and then write your own. Each file
runs a self-test when executed (\code{python3 code/ch04\_astar.py}), and the numbers
quoted in the worked examples were produced by that code. The figures that show
algorithm output were generated by the scripts in \code{code/figures/} from fixed
random seeds, so you can regenerate and modify them. The listings in the text are
excerpts; the files contain the complete programs.

\section*{Sources}

None of the algorithms in this book is new. Each chapter cites the paper in which its
algorithm was introduced and a textbook that treats it at greater length, and the
\emph{further reading} paragraph at the end of every chapter points to the
developments that the chapter leaves out. The bibliography lists all sources. Where
the book simplifies an algorithm for teaching, it says so and names the full version.

\section*{Errata}

The manuscript was reviewed in several rounds for technical accuracy, completeness
against the training plan, and clarity. Mistakes will nevertheless remain; if you find
one, the \LaTeX{} sources and code are distributed with the book precisely so that it
can be corrected.
